
考虑扰动\BookRecordIndex{T0209}柯西--黎曼方程
\begin{equation*}
  \tag{$*$}
  \bar\partial_J(u)=\nu(u),
\end{equation*}
的解往往很有用，其中 $\nu(u)\in\Omega^{0,1}(\Sigma,u^*TM)$。产生这种扰动的一种自然方式
如下。令 $\nu$ 为丛
\[
  \pi_1^*\Lambda^{0,1}T^*\Sigma\otimes\pi_2^*TM
  \longrightarrow \Sigma\times M
\]
的一个截面；它在 $(z,x)$ 处的纤维是复反线性同态 $T_z\Sigma\to T_xM$ 所成的空间。
然后，把 $\nu(u)\in\Omega^{0,1}(\Sigma,u^*TM)$ 定义为 $\nu$ 在嵌入
$\widetilde u:\Sigma\to\Sigma\times M$ 下的拉回，其中
\[
  \widetilde u(z):=(z,u(z)).
\]
在这种情形，\personlater{geluomofu}{格罗莫夫}{Gromov}在 \srcite{160} 中说明了如何在乘积
$\widetilde M:=\Sigma\times M$ 上定义\BookRecordIndex{T0088}殆复结构 $\widetilde J_\nu$，使得
$\widetilde M$ 的 $\widetilde J_\nu$-全纯截面恰好是方程 $(*)$ 的解的图像
$\widetilde u$。因此，我们可以用前面的方法分析 $(*)$ 的解。

\BookRecordIndex{T0370}虚模循环通常构造为：其元素是某个一般扰动 $\nu$ 所对应的 $(*)$ 的解。在这种情形，
人们只关心解的计数，所以解的精确性质并不重要。本章的重点略有不同。我们把注意力限制在
一类几何上自然的扰动上；对这类扰动，$(*)$ 是 弗洛尔方程的无坐标版本。此时解具有几何
意义；下一章将看到一些应用。

注意，每个扰动一形式 $\nu$ 在 $\{z\}\times M$ 上的限制，是一个从 $T_z\Sigma$ 到
$M$ 上向量场空间的线性映射之复反线性部分。一种格外良好的情形是：这个映射的像中的
向量场都是\BookRecordIndex{T0160}哈密顿向量场。此时称 $\nu$ 为\emph{\BookRecordIndex{T0157}哈密顿扰动}。我们将说明，相应的\BookRecordIndex{T0088}殆复结构
$\widetilde J_\nu$ 可用纤维丛 $\Sigma\times M\to\Sigma$ 上的联络作出自然的几何解释。

到这里已经没有必要局限于平凡丛；可以改为考虑任意\BookRecordIndex{T0200}局部哈密顿纤维化。于是，我们要考察
黎曼曲面上一般纤维丛 $\widetilde M\to\Sigma$ 的截面。这一框架有许多有趣的应用（见
\personlater{geluomofu}{格罗莫夫}{Gromov}\srcite{160}、Floer
\srcite{113,114,116}、Seidel \srcite{362,363}、Polterovich
\srcite{325,326,327}、Entov \srcite{104}、Akveld--Salamon \srcite{11}、
Piunikhin--Salamon--Schwarz \srcite{323}、McDuff \srcite{269}）。不过，本书许多结果
并不需要把理论推广到辛纤维丛；本章的写法允许读者略过这一推广。

第一节介绍带\BookRecordIndex{T0157}哈密顿扰动的\BookRecordIndex{T0209}柯西--黎曼方程，并解释它们与乘积丛上的辛联络之间的关系；
同时建立本章其余部分所用的记号。第~\ref{sec:8-2} 节把这些概念推广到更一般的纤维丛。
第~\ref{sec:8-3} 节讨论全纯截面的\BookRecordIndex{T0246}模空间，并概述相应\BookRecordIndex{T0166}横截性结果的证明。第~\ref{sec:8-4} 节考察
纤维中的全纯球面，第~\ref{sec:8-5} 节引入伪全纯截面的\BookRecordIndex{T0337}伪循环。第~\ref{sec:8-6} 节定义相应的
\BookRecordIndex{T0143}格罗莫夫--威滕不变量，并解释几个基本例子。整章可以看作第~\ref{sec:6-7} 节与
第~\ref{sec:7-3} 节在稍一般背景下的扩展版。
我们始终允许底流形 $\Sigma$ 有任意亏格。最终定义的\BookRecordIndex{T0143}格罗莫夫--威滕不变量计数
$J$-全纯截面 $u:\Sigma\to\widetilde M$ 的数目，其中底空间带有固定复结构 $j_\Sigma$。
最后一个例子将展示这一背景下出现的一些新特点。

\section{平凡丛}\label{sec:8-1}

本节讨论扰动\BookRecordIndex{T0209}柯西--黎曼方程 $(*)$ 的解，其中扰动
$\nu(u)\in\Omega^{0,1}(\Sigma,u^*TM)$ 是哈密顿的。这意味着，$\nu$ 是 $\Sigma$ 上一个
一形式的复反线性部分，而该一形式取值于 $M$ 上的\BookRecordIndex{T0160}哈密顿向量场空间。每个这样的扰动都由
一个一形式
\[
  H\in\Omega^1(\Sigma,C_0^\infty(M))
\]
唯一生成；它取值于 $M$ 上平均值为零的光滑函数空间 $C_0^\infty(M)$。我们把 $H$ 写成
\[
  T_z\Sigma\longrightarrow C_0^\infty(M):\quad
  \zeta\longmapsto H_\zeta.
\]
相应的、取值于\BookRecordIndex{T0160}哈密顿向量场空间的一形式为 $\zeta\mapsto X_{H_\zeta}$。给定光滑映射
$u:\Sigma\to M$，用 $X_H(u)\in\Omega^1(\Sigma,u^*TM)$ 表示沿 $u$ 的一形式
\[
  T_z\Sigma\longrightarrow T_{u(z)}M:\quad
  \zeta\longmapsto X_{H_\zeta}(u(z)),
\]
其值位于拉回切丛中。现在固定一族光滑的、被 $\omega$ \BookRecordIndex{T0372}驯服的\BookRecordIndex{T0088}殆复结构
\[
  \Sigma\longrightarrow\mathcal J_\tau(M,\omega):\quad z\longmapsto J_z.
\]
于是扰动\BookRecordIndex{T0209}柯西--黎曼方程 $(*)$ 具有形式
\begin{equation}\label{eq:8-1-1}
  \bar\partial_J(u)+X_H(u)^{0,1}=0,
\end{equation}
其中 $X_H(u)^{0,1}$ 表示 $X_H(u)$ 的复反线性部分：
\[
  X_H(u)^{0,1}(z)
  :=\frac12\bigl(X_H(u(z))+J_z\circ X_H(u(z))\circ j_\Sigma(z)\bigr).
\]
注意，上式第一项是\BookRecordIndex{T0160}哈密顿向量场，而第二项是梯度向量场。它们与
式~\eqref{eq:8-1-1} 左端的 $\bar\partial_J(u)$ 合在一起，构成\emph{协变导数}
\[
  d_H(u):=du+X_H(u)
\]
的复反线性部分，记作
\[
  \bar\partial_{J,H}(u):=d_H(u)^{0,1}
  =\bar\partial_J(u)+X_H(u)^{0,1}.
\]

\paragraph{\BookRecordIndex{T0155}哈密顿联络。}
为了从几何上理解式~\eqref{eq:8-1-1}，适宜在乘积流形
$\widetilde M:=\Sigma\times M$ 上工作，并用 $u$ 的图像
$\widetilde u:\Sigma\to\Sigma\times M$ 代替 $u$，其中
\[
  \widetilde u(z):=(z,u(z)).
\]
这样，$\widetilde u$ 是 $\widetilde M$ 的截面；相对于与扰动 $H$ 对应的辛联络，
它的协变导数 $d_H(u)$ 可以解释为底空间上的、沿 $\widetilde u$ 取值于水平切丛的一形式。
更明确地，每个 $H\in\Omega^1(\Sigma,C_0^\infty(M))$ 都给出
$\widetilde M$ 上的实值一形式 $\widetilde\sigma_H\in\Omega^1(\widetilde M)$，定义为
\[
  \widetilde\sigma_H(\zeta,\xi):=H_\zeta(x),
  \qquad \zeta\in T_z\Sigma,\quad \xi\in T_xM.
\]
注意，这不是 $H$ 到 $\widetilde M$ 的拉回（后者取值于 $C_0^\infty(M)$ 而非
$\mathbb R$），而是先拉回，再把所得函数在 $M$ 的相应点处求值。我们把
$\widetilde M$ 看成经由显然的投影
\[
  \pi:\widetilde M=\Sigma\times M\longrightarrow\Sigma
\]
\BookRecordIndex{T0279}纤维化在 $\Sigma$ 上。

一形式 $\widetilde\sigma_H$ 是水平的，即它在\BookRecordIndex{T0308}竖直切丛 $\ker d\pi$ 上消失。然而，
$d\widetilde\sigma_H$ 没有这一性质：它在一对竖直向量上消失，但只有一个向量竖直时完全
可能非零。考虑二形式
\begin{equation}\label{eq:8-1-2}
  \widetilde\omega=\widetilde\omega_H
  :=\operatorname{pr}^*\omega-d\widetilde\sigma_H
  \in\Omega^2(\widetilde M),
\end{equation}
其中 $\operatorname{pr}:\widetilde M\to M$ 是显然的投影。在 $\Sigma$ 的坐标片 $U$ 上取
局部\BookRecordIndex{T0149}共形坐标 $z=s+it\in\mathbb C$，则
$H\in\Omega^1(\Sigma,C_0^\infty(M))$ 可写成
\[
  H=F\,ds+G\,dt,
  \qquad F,G\in C^\infty(U,C_0^\infty(M)).
\]
因而
\begin{equation}\label{eq:8-1-3}
  \widetilde\omega_H
  =\omega-d'F\wedge ds-d'G\wedge dt
   +(\partial_tF-\partial_sG)\,ds\wedge dt,
\end{equation}
其中 $d'$ 表示对 $M$ 中变量取外微分。

二形式 $\widetilde\omega_H$ 是纤维上形式 $\omega$ 的一个闭延拓，因而在
Guillemin--Lerman--Sternberg \srcite{163} 和 McDuff--Salamon \srcite{277} 第 6 章的
术语中是一个\emph{\BookRecordIndex{T0224}联络形式}。相应的 $T\widetilde M$ 中的\BookRecordIndex{T0314}水平分布，在
$\widetilde M$ 的点 $\widetilde x$ 处为
\begin{equation}\label{eq:8-1-4}
  \operatorname{Hor}_{H;\widetilde x}
  :=\bigl\{\widetilde v\in T_{\widetilde x}\widetilde M
  \mid \widetilde\omega_H(\widetilde v,\widetilde w)=0
  \text{ 对所有 }\widetilde w\in\ker d\pi(\widetilde x)\bigr\}.
\end{equation}

尽管 $\widetilde\omega_H$ 从来不是 $\widetilde M$ 上的\BookRecordIndex{T0363}辛形式（见习题~\ref{ex:8-1-3}），
它在每条纤维上的限制都是辛的；当 $M$ 紧时，给它加上底空间 $\Sigma$ 上某个面积形式拉回
的适当大倍数，就能使它成为\BookRecordIndex{T0363}辛形式。因此，对 $\Sigma$ 中任意嵌入弧 $\gamma$，
$\widetilde\omega_H$ 在 $\widetilde M$ 中超曲面 $\pi^{-1}(\gamma)$ 上的限制都具有最大秩
\sourcefootnote{原书此处印作 ``nondegenerate''；但 $\pi^{-1}(\gamma)$ 是奇维超曲面，
而且紧接着的句子明确给出这一限制的一维核，故应理解为具有最大秩。}，其核恰由
$\gamma$ 的切向量的\BookRecordIndex{T0317}水平提升张成。因此，沿 $\gamma$ 的\BookRecordIndex{T0263}平行移动由这个超曲面上特征流的适当参数化给出，从而
保持纤维上的\BookRecordIndex{T0363}辛形式。如下引理表明，它实际上由 $M$ 的一族哈密顿辛同胚给出。

\begin{lemma}\label{lem:8-1-1}
设 $H\in\Omega^1(\Sigma,C_0^\infty(M))$。相对于
$\widetilde\omega_H$ 所确定的联络~\eqref{eq:8-1-4}，沿 $\Sigma$ 中任意路径的\BookRecordIndex{T0263}平行移动，
都是 $M$ 的哈密顿辛同胚。
\end{lemma}

\begin{proof}
令 $\gamma:[0,1]\to\Sigma$ 为光滑路径，$x:[0,1]\to M$ 为其\BookRecordIndex{T0317}水平提升。这意味着
\begin{align*}
  0
  &=\widetilde\omega_H\bigl((\dot\gamma(t),\dot x(t)),(0,\xi)\bigr)\\
  &=\omega(\dot x,\xi)
    -d\widetilde\sigma_H\bigl((\dot\gamma,\dot x),(0,\xi)\bigr)\\
  &=\omega(\dot x,\xi)+dH_{\dot\gamma}(x)\xi\\
  &=\omega\bigl(\dot x+X_{H_{\dot\gamma}}(x),\xi\bigr)
\end{align*}
对每个 $t\in[0,1]$ 和每个 $\xi\in T_{x(t)}M$ 都成立。等价地，$x$ 满足含时哈密顿微分方程
\[
  \dot x+X_{H_{\dot\gamma}}(x)=0.
\]
沿 $\gamma$ 的\BookRecordIndex{T0263}平行移动就是 $x(0)\mapsto x(1)$ 的时间 $1$ 映射，引理得证。
\end{proof}

\paragraph{曲率。}
$\Sigma\times M$ 上的\BookRecordIndex{T0155}哈密顿联络，是平凡 $G$-丛 $P=\Sigma\times G$ 上联络的特殊情形。
这里结构群 $G$ 是 $M$ 的哈密顿辛同胚群，其李代数
$\mathfrak g=\operatorname{Lie}(G)$ 是平均值为零的哈密顿函数空间 $C_0^\infty(M)$，
李代数结构由 Poisson 括号给出。\BookRecordIndex{T0223}联络势是 $\Sigma$ 上取值于李代数的一形式 $H$，其曲率
是取值于李代数的二形式。把曲率形式写成 $R_H\,\mathrm{dvol}_\Sigma$ 很方便，其中
$R_H:\Sigma\times M\to\mathbb R$ 是光滑函数。它由
\[
  R_H\,\mathrm{dvol}_\Sigma
  =dH+\frac12\{H\wedge H\}
  \in\Omega^2(\Sigma,C_0^\infty(M))
\]
给出。从几何上说，对任意一对向量场
$v,w\in\operatorname{Vect}(\Sigma)$，曲率在 $(v,w)$ 上的值，是它们的\BookRecordIndex{T0317}水平提升
$v^\#,w^\#\in\operatorname{Vect}(\widetilde M)$ 的李括号 $[v^\#,w^\#]$ 的竖直部分。
这个竖直部分在纤维 $\{z\}\times M$ 上的限制，是哈密顿函数
\[
  R_H(z,\mathord\cdot)\,\mathrm{dvol}_\Sigma(v(z),w(z))
  \in C_0^\infty(M)
\]
所对应的 $M$ 上的\BookRecordIndex{T0160}哈密顿向量场。

\begin{exercise}\label{ex:8-1-2}
设 $\widetilde\omega_H$ 在局部坐标中由式~\eqref{eq:8-1-3} 给出。证明曲率形式为
\[
  R_H\lambda=\partial_sG-\partial_tF+\{F,G\},
\]
其中 $\lambda\,ds\wedge dt$ 是 $U$ 上的体积形式，$\{\mathord\cdot,\mathord\cdot\}$
表示 $M$ 上的 Poisson 括号。由此\BookRecordIndex{T0326}推出，对 $\Sigma$ 上任意一对向量场 $v,w$，都有
\[
  R_H\,\mathrm{dvol}(v,w)=-\widetilde\omega_H(v^\#,w^\#).
\]
提示：利用向量场
\[
  \partial_s^\#=(\partial_s,-X_F),
  \qquad
  \partial_t^\#=(\partial_t,-X_G)
\]
是水平的这一事实，并计算 $[\partial_s^\#,\partial_t^\#]$ 的竖直部分。
\end{exercise}

\begin{exercise}\label{ex:8-1-3}
利用习题~\ref{ex:8-1-2} 的结果证明，曲率 $R_H:\Sigma\times M\to\mathbb R$ 由公式
\[
  \frac{\widetilde\omega_H^{n+1}}{(n+1)!}
  =-R_H\,\pi^*\mathrm{dvol}_\Sigma\wedge
    \operatorname{pr}^*\frac{\omega^n}{n!}
\]
确定，其中 $\pi:\Sigma\times M\to\Sigma$ 与
$\operatorname{pr}:\Sigma\times M\to M$ 是显然的投影。因为 $R_H$ 在每条纤维上的平均值
为零，这个公式说明二形式 $\widetilde\omega_H\in\Omega^2(\widetilde M)$ 不可能是\BookRecordIndex{T0363}辛形式。
不过，若光滑函数 $\kappa:\Sigma\to\mathbb R$ 满足
\begin{equation}\label{eq:8-1-5}
  \kappa(z)>R_H(z,x)
\end{equation}
对所有 $z\in\Sigma$ 和 $x\in M$ 都成立，则二形式
\[
  \widetilde\omega_{H,\kappa}
  :=\widetilde\omega_H+\pi^*(\kappa\,\mathrm{dvol}_\Sigma)
\]
是 $\widetilde M$ 上的\BookRecordIndex{T0363}辛形式。
\end{exercise}

\paragraph{\BookRecordIndex{T0088}殆复结构。}
现在介绍 \personlater{geluomofu}{格罗莫夫}{Gromov}的一个构造，它表明，扰动\BookRecordIndex{T0209}柯西--黎曼
方程的解的图像，可以解释为 $\widetilde M$ 上关于某个适当\BookRecordIndex{T0088}殆复结构的全纯曲线。
相关的\BookRecordIndex{T0088}殆复结构 $\widetilde J_H$ 具有如下性质：它在纤维 $\{z\}\times M$ 上的限制等于
$J_z$；投影 $\pi$ 关于 $\widetilde J_H$ 与 $j_\Sigma$ 是全纯的；而且
$\widetilde J_H$ 适配于\BookRecordIndex{T0224}联络形式 $\widetilde\omega_H$，即它保持\BookRecordIndex{T0314}水平分布
$\operatorname{Hor}_H$。给定 $\{J_z\}$ 与 $j_\Sigma$ 后，这些条件唯一确定
$\widetilde J_H$。它由
\[
  \widetilde J_H(z,x):=
  \begin{pmatrix}
    j_\Sigma(z) & 0\\
    J(z,x)\circ X_H(z,x)-X_H(z,x)\circ j_\Sigma(z) & J(z,x)
  \end{pmatrix}
\]
给出；这里 $X_H(z,x):T_z\Sigma\to T_xM$ 表示线性映射
$\zeta\mapsto X_{H_\zeta}(x)$，并记 $J(z,x):=J_z(x)$。

\begin{exercise}\label{ex:8-1-4}
令 $\kappa:\Sigma\to\mathbb R$ 为满足式~\eqref{eq:8-1-5} 的光滑函数。证明：
$\widetilde J_H$ 被\BookRecordIndex{T0363}辛形式 $\widetilde\omega_{H,\kappa}$ \BookRecordIndex{T0372}驯服，当且仅当每个 $J_z$ 都被
$\omega$ \BookRecordIndex{T0372}驯服。类似地，$\widetilde J_H$ 与 $\widetilde\omega_{H,\kappa}$ \BookRecordIndex{T0348}相容，当且仅当
每个 $J_z$ 都与 $\omega$ \BookRecordIndex{T0348}相容。提示：使用公式
\begin{align*}
  &\widetilde\omega_H
    \bigl((\sigma,\tau,\xi),(\sigma',\tau',\xi')\bigr)\\
  &\quad=
  \omega\bigl(\xi+\sigma X_F+\tau X_G,
               \xi'+\sigma'X_F+\tau'X_G\bigr)
  -(\sigma\tau'-\tau\sigma')R_H,
\end{align*}
其中 $\sigma,\tau,\sigma',\tau'\in\mathbb R$，$\xi,\xi'\in T_xM$。
\end{exercise}

\begin{exercise}\label{ex:8-1-5}
证明：光滑映射 $u:\Sigma\to M$ 是式~\eqref{eq:8-1-1} 的解，当且仅当其图像
$\widetilde u:\Sigma\to\widetilde M$，$\widetilde u(z):=(z,u(z))$，是
$\widetilde M=\Sigma\times M$ 的 $\widetilde J_H$-全纯截面。
\end{exercise}

\paragraph{局部坐标。}
在 $\Sigma$ 的一个开子集上选取\BookRecordIndex{T0149}共形坐标 $z=s+it\in U\subset\mathbb C$。在这样的坐标中，
一形式 $H\in\Omega^1(\Sigma,C_0^\infty(M))$ 可写成
$H=F\,ds+G\,dt$，其中 $F,G\in C^\infty(U,C_0^\infty(M))$。如上所见，
\[
  \widetilde\omega_H
  =\omega-d'F\wedge ds-d'G\wedge dt
   +(\partial_tF-\partial_sG)\,ds\wedge dt,
\]
其中 $d'$ 表示对 $M$ 中变量取外微分，而曲率的局部公式为
\[
  R_H\lambda=\partial_sG-\partial_tF+\{F,G\},
\]
这里 $\lambda\,ds\wedge dt$ 表示 $U$ 上的体积形式。在 $TU$ 的\BookRecordIndex{T0067}标架
$\partial_s,\partial_t$ 下，水平子空间由向量 $(1,0,-X_F)$ 与 $(0,1,-X_G)$ 张成。因此
\begin{equation}\label{eq:8-1-6}
  \widetilde J_H=
  \begin{pmatrix}
    0&-1&0\\
    1&0&0\\
    JX_F-X_G&JX_G+X_F&J_{s,t}
  \end{pmatrix}.
\end{equation}

在局部坐标中，协变导数 $d_H(u)$ 具有形式
\[
  d_H(u)=(\partial_su+X_F(u))\,ds+(\partial_tu+X_G(u))\,dt,
\]
其中 $X_F=X_{F_{s,t}}$ 与 $X_G=X_{G_{s,t}}$ 分别表示 $F$ 与 $G$ 的\BookRecordIndex{T0160}哈密顿向量场。
因此，光滑映射 $u:\Sigma\to M$ 满足式~\eqref{eq:8-1-1}，当且仅当它在局部坐标中满足
\begin{equation}\label{eq:8-1-7}
  \partial_su+X_F(s,t,u)
  +J(s,t,u)\bigl(\partial_tu+X_G(s,t,u)\bigr)=0.
\end{equation}
读者可以验证，$u:U\to M$ 满足这个方程，当且仅当映射
$(s,t)\mapsto(s,t,u(s,t))$ 是 $\widetilde J_H$-全纯曲线。

\paragraph{\BookRecordIndex{T0249}能量恒等式。}
固定一个\BookRecordIndex{T0157}哈密顿扰动 $H\in\Omega^1(\Sigma,C_0^\infty(M))$，以及一族光滑的、被
$\omega$ \BookRecordIndex{T0372}驯服的\BookRecordIndex{T0088}殆复结构
\[
  J=\{J_z\}_{z\in\Sigma}
  \in C^\infty(\Sigma,\mathcal J_\tau(M,\omega)).
\]
光滑映射 $u:\Sigma\to M$ 的能量，定义为其协变导数 $d_H(u)$ 的 $L^2$-范数平方的一半：
\begin{equation}\label{eq:8-1-8}
  E_H(u):=\frac12\int_\Sigma\lvert d_H(u)\rvert_J^2\,
  \mathrm{dvol}_\Sigma.
\end{equation}
当 $H=0$ 时，这就归结为通常的能量概念。但现在非常值映射的能量也可能为零。事实上，
$E_H(u)=0$ 当且仅当 $\widetilde u(z)=(z,u(z))$ 是 $\widetilde M$ 的\BookRecordIndex{T0315}水平截面。\BookRecordIndex{T0249}能量恒等式
涉及曲率 $R_H$，具有如下形式。

\begin{lemma}\label{lem:8-1-6}
设 $u:\Sigma\to M$ 是式~\eqref{eq:8-1-1} 的解。则
\begin{equation}\label{eq:8-1-9}
  E_H(u)=\int_\Sigma u^*\omega
  +\int_\Sigma R_H(\widetilde u)\,\mathrm{dvol}_\Sigma.
\end{equation}
\end{lemma}

\begin{proof}
在 $U$ 上的局部坐标中，被积函数具有形式
\begin{align*}
  \lvert\partial_su+X_F(u)\rvert^2
  &:=\omega\bigl(\partial_su+X_F(u),\partial_tu+X_G(u)\bigr)\\
  &=\omega(\partial_su,\partial_tu)
    +d'F(u)\partial_tu-d'G(u)\partial_su+\{F,G\}\circ u\\
  &=\widetilde\omega_H\bigl((1,0,\partial_su),(0,1,\partial_tu)\bigr)
    +(\partial_sG-\partial_tF+\{F,G\})\circ u\\
  &=\widetilde\omega_H(\partial_s\widetilde u,\partial_t\widetilde u)
    +\frac{R_H\circ\widetilde u}{\lambda},
\end{align*}
其中 $\lambda\,ds\wedge dt$ 是 $U$ 上的体积形式。因为 $\widetilde\omega_H$ 与
$\operatorname{pr}^*\omega$ 相差一个\BookRecordIndex{T0403}正合形式，$\widetilde u^*\widetilde\omega_H$ 在
$\Sigma$ 上的积分等于 $u^*\omega$ 的积分。这就证明了式~\eqref{eq:8-1-9}。
\end{proof}

\begin{remark}\label{rem:8-1-7}
曲率 $R=R_H$ 的 \emph{\BookRecordIndex{T0028}Hofer 范数}定义为
\[
  \lVert R_H\rVert
  :=\int_\Sigma
  \left(\max_{x\in M}R_H(z,x)-\min_{x\in M}R_H(z,x)\right)
  \mathrm{dvol}_\Sigma(z).
\]
注意，这个范数与 $\Sigma$ 上体积形式的选取无关。由\BookRecordIndex{T0249}能量恒等式~\eqref{eq:8-1-9}，
式~\eqref{eq:8-1-1} 的每个解都满足
\[
  E_H(u)\le\int_\Sigma u^*\omega+\lVert R_H\rVert.
\]
这个能量上界不再是纯拓扑量，因为它还依赖于联络的曲率。它的优点是：即使图像
$\widetilde u(z)=(z,u(z))$ 的总能量没有定义或无界，它仍然有定义且有界；例如，注
~\ref{rem:8-1-9} 所考虑的 弗洛尔方程就有非紧定义域。
\end{remark}

\begin{exercise}\label{ex:8-1-8}
令 $u$ 为式~\eqref{eq:8-1-1} 的解，并考虑其图像
$\widetilde u:\Sigma\to\widetilde M$。以
$E_{H,\kappa}(\widetilde u)$ 表示它关于 $\widetilde\omega_{H,\kappa}$ 与
$\widetilde J_H$ 所确定度量的能量，其中 $\kappa:\Sigma\to\mathbb R$ 满足
式~\eqref{eq:8-1-5}（见习题~\ref{ex:8-1-3}、\ref{ex:8-1-4}、\ref{ex:8-1-5}）。证明
\begin{align*}
  E_{H,\kappa}(\widetilde u)
  &=\int_\Sigma
    \left(\frac12\lvert d_Hu\rvert_J^2+\kappa-R_H(\widetilde u)\right)
    \mathrm{dvol}_\Sigma\\
  &=\int_\Sigma\widetilde u^*\widetilde\omega_{H,\kappa}
   =\int_\Sigma u^*\omega+\int_\Sigma\kappa\,\mathrm{dvol}_\Sigma.
\end{align*}
由此得到估计
\[
  E_H(u)\le\int_\Sigma u^*\omega+\int_\Sigma\kappa\,\mathrm{dvol}_\Sigma.
\]
\end{exercise}

\begin{remark}[弗洛尔方程]\label{rem:8-1-9}
考虑带标准复结构的圆柱 $\Sigma:=\mathbb R\times S^1$ 上的扰动\BookRecordIndex{T0209}柯西--黎曼
方程~\eqref{eq:8-1-7}。此时有整体全纯坐标 $s+it$，其中
$s\in\mathbb R$，$t\in S^1=\mathbb R/\mathbb Z$。假设\BookRecordIndex{T0088}殆复结构和\BookRecordIndex{T0157}哈密顿扰动具有与
$\mathbb R$-作用相联系的额外对称性：$J_{s,t}=J_t$ 与变量 $s$ 无关，而\BookRecordIndex{T0157}哈密顿扰动具有
形式 $F\,ds+G\,dt$，其中对所有 $s,t$ 都有 $F_{s,t}=0$、$G_{s,t}=-H_t$。于是
式~\eqref{eq:8-1-7} 化为 Floer 所研究的方程
\[
  \partial_su+J_t(u)\bigl(\partial_tu-X_{H_t}(u)\bigr)=0.
\]
因此，我们的扰动\BookRecordIndex{T0209}柯西--黎曼方程可解释为依赖于 $s$ 的 弗洛尔方程。第~\ref{ch:applications-symplectic-topology} 章的一些应用
将使用这一观察。
\end{remark}

\begin{remark}[拉格朗日边界条件]\label{rem:8-1-10}
上面对 $J$-全纯截面的讨论可推广到非紧纤维 $(M,\omega)$，也可推广到带边界
$\partial\Sigma$ 的紧黎曼曲面 $\Sigma$。此时\BookRecordIndex{T0249}能量恒等式会多出一个边界项，它来自
$\widetilde u^*d\widetilde\sigma_H$ 的积分。现在\BookRecordIndex{T0249}能量恒等式为
\begin{equation}\label{eq:8-1-10}
  E_H(u)=\int_\Sigma u^*\omega
  +\int_\Sigma R_H(u)\,\mathrm{dvol}_\Sigma
  -\int_{\partial\Sigma}H(u),
\end{equation}
对式~\eqref{eq:8-1-1} 的每个解 $u:\Sigma\to M$ 都成立；这里
$H(u)\in\Omega^1(\Sigma)$ 表示一形式
\[
  T_z\Sigma\longrightarrow\mathbb R:\quad
  \zeta\longmapsto H_\zeta(u(z)).
\]
在有边界的情形，考察拉格朗日子丛
$\widetilde L\subset\partial\Sigma\times M$ 很有意思；其纤维
\[
  L_z:=\{x\in M\mid(z,x)\in\widetilde L\}
\]
对每个 $z\in\partial\Sigma$ 都是 $M$ 的\BookRecordIndex{T0219}拉格朗日子流形。在这种情形，形式
$\widetilde\omega_H=\widetilde\omega-d\widetilde\sigma_H$ 在
$T\widetilde L$ 上消失，当且仅当 $\widetilde L$ 在沿 $\partial\Sigma$ 的\BookRecordIndex{T0263}平行移动下不变
（见 \srcite{11}，引理 3.1）。若此条件成立，就可以考虑满足边界条件
\begin{equation}\label{eq:8-1-11}
  u(z)\in L_z\qquad(z\in\partial\Sigma)
\end{equation}
的式~\eqref{eq:8-1-1} 的解。Akveld--Salamon 在 \srcite{11} 中用
式~\eqref{eq:8-1-1} 与~\eqref{eq:8-1-11} 的解，从下方估计固定\BookRecordIndex{T0158}哈密顿同痕类中\BookRecordIndex{T0219}拉格朗日
子流形环路的 Hofer 长度。
\end{remark}

在应用中，我们常假设\BookRecordIndex{T0088}殆复结构 $J=\{J_z\}$ 固定，而扰动 $H$ 在集合
\[
  \mathcal H(\kappa)\subset\Omega^1(\Sigma,C_0^\infty(M))
\]
中变化；这个集合由曲率被 $\kappa$ 控制的所有一形式组成。由此得到
$\widetilde M$ 上一族\BookRecordIndex{T0088}殆复结构 $\widetilde J_H$。第~\ref{sec:8-3} 节将看到，这个族足够大，
通常的\BookRecordIndex{T0166}横截性论证可以进行。换言之，对一般的 $H\in\mathcal H(\kappa)$，表示同调类
$A\in H_2(M;\mathbb Z)$ 的式~\eqref{eq:8-1-1} 的解所成的\BookRecordIndex{T0246}模空间
$\mathcal M(A;J,H)$ 是光滑流形。于是第~\ref{sec:6-7} 节的论证适用于当前情形，从而可以用扰动
\BookRecordIndex{T0209}柯西--黎曼方程的解定义\BookRecordIndex{T0143}格罗莫夫--威滕不变量。本章稍后会在\BookRecordIndex{T0200}局部哈密顿纤维化
$\pi:\widetilde M\to\Sigma$ 的更一般情形下给出细节。第~\ref{ch:applications-symplectic-topology} 章的许多应用不需要这一推广，
读者可以直接跳到第~\ref{sec:8-3} 节。
